% ---------------------------------------------------------------------------
% seed_522.tex --- the initial input to the workflow
%
% This worked ((5,2,2)) derivation was written by hand by one of the authors
% (B.Z.) and given to GPT-5 as the starting point of the project.  At that stage
% there was no agent framework: a human posed this document directly to the
% model.  It is reproduced in full as Sec. 2.3 of arXiv:2510.20728v1
% (23 October 2025), and it is the same code that Examples/D2/Ex522.lean
% certifies -- same supports, same rational probabilities, same target
% Z-expectations.
%
% Kept verbatim.  Nothing here has been edited for the repository.
% ---------------------------------------------------------------------------

\documentclass{article}
\usepackage{amsmath,amssymb}
\usepackage{braket}

\begin{document}

\section*{Subset-Sum Quantum Code: Detailed Summary}

\subsection*{Step 1: Weight Vector and Modulus}

\begin{itemize}
  \item Number of qubits: \( n = 5 \)
  \item Weights: \( \vec{w} = (1,1,2,2,2) \mod 7 \)
  \item Modulus: \( m = 7 \)
  \item Target sums: \( S_0 = 0 \), \( S_1 = 4 \)
\end{itemize}

\subsection*{Step 2: Codewords}

Codewords are defined by modular subset sums:
\[
C_0 = \left\{ x \in \{0,1\}^5 : \sum_i w_i x_i \equiv 0 \mod 7 \right\}
\]
\[
C_1 = \left\{ x \in \{0,1\}^5 : \sum_i w_i x_i \equiv 4 \mod 7 \right\}
\]

Explicitly:
\[
C_0 = \{\ket{00000},\ \ket{01111},\ \ket{10111} \}
\]
\[
C_1 = \{\ket{00011},\ \ket{00101},\ \ket{00110},\ \ket{11001},\ \ket{11010},\ \ket{11100} \}
\]

\subsection*{Step 3: Hamming Distance}

Minimum Hamming distance between any pair of codewords in \( C = C_0 \cup C_1 \):
\[
d(C) = 2
\]

\subsection*{Step 4: Logical States (Parameterized Form)}

Let:
\[
\ket{0_L} = x_0 \ket{00000} + x_1 \ket{01111} + x_2 \ket{10111}
\]
\[
\ket{1_L} = y_0 \ket{00011} + y_1 \ket{00101} + y_2 \ket{00110} + y_3 \ket{11001} + y_4 \ket{11010} + y_5 \ket{11100}
\]
where the coefficients \( x_i \) and \( y_j \) are complex amplitudes subject to normalization and the Knill--Laflamme conditions.

\subsection*{Step 5: Action of Transversal Gate}

Apply
\[
U = Z_1\left(\tfrac{2\pi}{7}\right) Z_2\left(\tfrac{2\pi}{7}\right) Z_3\left(\tfrac{4\pi}{7}\right) Z_4\left(\tfrac{4\pi}{7}\right) Z_5\left(\tfrac{4\pi}{7}\right)
\]

Then:
\[
U \ket{0_L} = \ket{0_L}, \quad U \ket{1_L} = \omega_7^4 \ket{1_L}, \quad \omega_7 = e^{2\pi i/7}
\]

Logical gate:
\[
\overline{Z} = \text{diag}(1,\ \omega_7^4)
\]

\subsection*{Step 6: KL Conditions from \( Z_i \)}

Let \( X_j = |x_j|^2 \) and \( Y_k = |y_k|^2 \). Then the KL conditions require:

\[
\langle 0_L | Z_i | 0_L \rangle = \langle 1_L | Z_i | 1_L \rangle \quad \text{for } i = 1,\dots,5
\]

This gives the system:
\begin{align*}
X_0 + X_1 - X_2 &= Y_0 + Y_1 + Y_2 - Y_3 - Y_4 - Y_5 \\
X_0 - X_1 - X_2 &= Y_0 - Y_1 - Y_2 + Y_3 + Y_4 - Y_5 \\
X_0 - X_1 - X_2 &= -Y_0 + Y_1 - Y_2 + Y_3 - Y_4 + Y_5 \\
X_0 - X_1 - X_2 &= -Y_0 - Y_1 + Y_2 - Y_3 + Y_4 + Y_5 \\
X_0 - X_1 + X_2 &= Y_0 + Y_1 + Y_2 - Y_3 - Y_4 - Y_5
\end{align*}

\subsection*{Step 7: Normalization Conditions}

\[
X_0 + X_1 + X_2 = 1, \quad Y_0 + Y_1 + Y_2 + Y_3 + Y_4 + Y_5 = 1
\]

\subsection*{Step 8: General Solution (Gaussian Elimination)}

\[
\begin{aligned}
X_0 &= \frac{3}{7},\quad X_1 = X_2 = \frac{2}{7} \\
Y_0 &= \frac{1}{7} + Y_5,\quad Y_1 = \frac{1}{7} + Y_4 \\
Y_2 &= \frac{3}{7} - Y_4 - Y_5,\quad Y_3 = \frac{2}{7} - Y_4 - Y_5
\end{aligned}
\]
with \( Y_4, Y_5 \geq 0 \), \( Y_4 + Y_5 \leq \frac{2}{7} \)

\subsection*{Step 9: Special Solution with \( Y_4 = Y_5 = 0 \)}

\[
\begin{aligned}
\ket{0_L} &= \sqrt{\tfrac{3}{7}} \ket{00000} + \sqrt{\tfrac{2}{7}} \ket{01111} + \sqrt{\tfrac{2}{7}} \ket{10111} \\
\ket{1_L} &= \sqrt{\tfrac{1}{7}} \ket{00011} + \sqrt{\tfrac{1}{7}} \ket{00101} + \sqrt{\tfrac{3}{7}} \ket{00110} + \sqrt{\tfrac{2}{7}} \ket{11001}
\end{aligned}
\]

\subsection*{Step 10: KL Expectations}

\[
\langle Z_i \rangle =
\begin{cases}
\frac{3}{7}, & i = 1,2 \\
-\frac{1}{7}, & i = 3,4,5
\end{cases}
\quad \text{(holds for both \( \ket{0_L} \) and \( \ket{1_L} \))}
\]

\subsection*{Conclusion}

This construction yields a \textbf{\(((5,2,2))\) quantum code} that:
\begin{itemize}
  \item Satisfies the Knill–Laflamme error correction conditions for all single-qubit \(Z_i\) operators,
  \item Has minimum Hamming distance \(d = 2\),
  \item Supports a transversal implementation of the logical diagonal gate:
  \[
  \overline{Z} = \mathrm{diag}(1,\ \omega_7^4), \quad \omega_7 = e^{2\pi i/7}
  \]
  \item Encodes a 2-dimensional logical space (1 qubit) into 5 physical qubits.
\end{itemize}

\end{document}
